\documentclass{beamer}

% Кодировка и языки
\usepackage[T2A]{fontenc}        
\usepackage[utf8]{inputenc}   
\usepackage[russian]{babel}

% Поля
%\usepackage[left=25mm, right=15mm, top=20mm, bottom=20mm]{geometry} % Настройка полей

% Для вставки изображений
\usepackage{graphicx}

% Математические символы
\usepackage{amsmath}
\usepackage{amssymb} 
\usepackage{multirow}

\title{Домашнее задание по математическому анализу}
\author{Самир}
\date{\today}

\usetheme{Madrid}

\begin{document}

\begin{frame}
\maketitle
\end{frame}

\begin{frame}{Содержание работы}
В данной презентации рассмотрены следующие темы:
\begin{itemize}
    \item Определение и примеры замечательных пределов.
    \item Основные правила дифференцирования и таблица производных.
    \item Формула Тейлора и остаточный член в форме Лагранжа.
    \item Список используемой учебной литературы.
\end{itemize}
\end{frame}

\begin{frame}{Замечательные пределы}
В ходе выполнения домашнего задания были рассмотрены ключевые пределы:

\begin{block}{Первый замечательный предел}
\begin{equation}
\lim_{x\to0} \frac{\sin x}{x}=1
\label{eq:slide_sin_limit}  
\end{equation}
\end{block}

\begin{block}{Второй замечательный предел}
\begin{equation}
\lim_{x\to0} (1+x)^\frac{1}{x}=e   
\end{equation}
\end{block}
\end{frame}


\begin{frame}{Правила дифференцирования}
Основные правила вычисления производных для базовых операций:
\begin{enumerate} 
    \item \textbf{Линейность операции:} 
    $$(af + bg)' = af' + bg'$$
    \item \textbf{Производная произведения:} 
    $$(fg)' = f'g + fg'$$
    \item \textbf{Производная частного:} 
    $$\left(\frac{f}{g}\right)' = \frac{f'g - fg'}{g^2}$$
\end{enumerate} 
\end{frame}



\begin{frame}{Таблица производных}
Сводная таблица производных

\begin{table}
\centering
\small
\begin{tabular}{|c|c|c|}
\hline
\textbf{Класс функции} & \textbf{Функция $f(x)$} & \textbf{Производная $f'(x)$} \\
\hline
\multirow{2}{*}{Степенная} & $x^n$ & $n x^{n-1}$ \\
\cline{2-3}
 & $C$ (const) & $0$ \\
\hline
\multirow{2}{*}{Тригонометрическая} & $\sin x$ & $\cos x$ \\
\cline{2-3}
 & $\cos x$ & $-\sin x$ \\
\hline
Показательная & $e^x$ & $e^x$ \\
\hline
\end{tabular}
\caption{Основные производные элементачных функций}
\end{table}
\end{frame}


\begin{frame}{Формула Тейлора}
Разложение функции в окрестности точки $a$:
\begin{equation*}
f(x) = f(a) + f'(a)(x-a) + \frac{f''(a)}{2!}(x-a)^2 + \dots + \frac{f^{(n)}(a)}{n!}(x-a)^n + R_n
\end{equation*}

\begin{block}{Остаточный член в форме Лагранжа}
\begin{equation*}
R_n = \frac{f^{(n+1)}(\xi)}{(n+1)!}(x-a)^{n+1}
\end{equation*}
\end{block}
\end{frame}


\begin{frame}{Список литературы}
При подготовке материалов использовалась следующая учебная литература:

\begin{thebibliography}{1}
\bibitem{zorich_slide}
В.А.Зорич.
\textit{Математический анализ. Часть 1}.
\end{thebibliography}
\end{frame}

\end{document}